% GENERATED FILE. Edit canonical lessons, then run npm run generate:books.
% canonical-source-hash: fnv1a64:e84a3efc9bc44d4f
% canonical-lessons: MW-C224-phirno, MW-C224-kathputli, MW-C224-ladu, MW-C224-uchhalno, MW-C224-nagaro

\chapter{Fun at the Fair}
\label{ch:mw-fun-at-the-fair}
\hlchaptermodality{224}

\begin{chapteropening}
\textbf{By the end of this chapter:} \emph{I can say what I will do and eat at the fair.}

Complete the last lesson of chapter 224: I can say what I will do and eat at the fair.
\end{chapteropening}

\section[\texorpdfstring{\mw{फिरणो} (phirṇo)}{फिरणो (phirṇo)}]{\mw{फिरणो} (phirṇo) — to wander about; to roam}

\label{lesson:MW-C224-phirno}



\begin{quote}
Before the new one: say the Marwadi for a flag, then the Marwadi for a show.
\end{quote}

\subsection*{You'll want to know: \mw{फिरणो}}
\textbf{\mw{फिरणो}} — \emph{phirṇo} — \textquotedblleft{}to wander about; to roam\textquotedblright{}.

\begin{quote}
\mw{हूं} \mw{काल} \mw{बाजार} \mw{मांय} \mw{फिरसूं।}
\end{quote}

\begin{quote}
\emph{hū̃ kāl bājār mā̃y phirsū̃.}
\end{quote}

\begin{quote}
Tomorrow I'll wander round the market.
\end{quote}

Past (he …): \mw{वो} \mw{बाजार} \mw{मांय} \mw{फिर्यो}, \emph{vo bājār mā̃y phiryo}.

\begin{cousinweb}
\textbf{\mw{फिरणो}} (\emph{phirṇo}) is inherited from a reconstructed \emph{*phirati}, \textquotedblleft{}moves, wanders, turns\textquotedblright{}; CDIAL lists the Old Marwari verb \emph{phiraï}, \textquotedblleft{}turns, walks about\textquotedblright{} (CDIAL 9078). \textbf{\mw{फेर}}, again, is from the same family.
\end{cousinweb}

\begin{grammarlens}[title={What you've built}]
One more word for saying what comes next.
\end{grammarlens}

\subsection*{Your turn}
\begin{itemize}
\raggedright
  \item \emph{Say it:} \emph{phirṇo}
  \item \emph{Say it:} \emph{phirṇo}, once more
  \item \emph{Say it:} \emph{hū̃ kāl bājār mā̃y phirsū̃.} — Tomorrow I'll wander round the market.
  \item \emph{Recall:} say the Marwadi for a flag, then the Marwadi for a show, then say \emph{phirṇo} again
\end{itemize}

\begin{tcolorbox}[breakable,colback=teal!4,colframe=teal!35!black,title={Before you move on}]
What does \mw{फिरणो} mean? (\textquotedblleft{}To wander about; to roam\textquotedblright{}.) Say it once more.
\end{tcolorbox}
\section[\texorpdfstring{\mw{कठपूतळी} (kaṭhpūtḷī)}{कठपूतळी (kaṭhpūtḷī)}]{\mw{कठपूतळी} (kaṭhpūtḷī) — a puppet}

\label{lesson:MW-C224-kathputli}



\begin{quote}
Before the new one: say the Marwadi for a show, then the Marwadi for to wander about.
\end{quote}

\subsection*{You'll want to know: \mw{कठपूतळी}}
\textbf{\mw{कठपूतळी}} — \emph{kaṭhpūtḷī} — \textquotedblleft{}a puppet\textquotedblright{}.

\begin{quote}
\mw{कठपूतळी} \mw{नाचै} \mw{है।}
\end{quote}

\begin{quote}
\emph{kaṭhpūtḷī nāchai hai.}
\end{quote}

\begin{quote}
The puppet dances.
\end{quote}

A feminine noun.

\begin{cousinweb}
\textbf{\mw{कठपूतळी}} (\emph{kaṭhpūtḷī}) is built from \emph{kāṭh}, \textquotedblleft{}wood\textquotedblright{} (CDIAL 3120), and \emph{pūtḷī}, \textquotedblleft{}a doll\textquotedblright{}, from a reconstructed \emph{*putrala}, \textquotedblleft{}a puppet\textquotedblright{}; CDIAL lists Old Marwari \emph{pūtaḷī}, \textquotedblleft{}an image\textquotedblright{} (CDIAL 8269). It is a wooden doll worked on strings.
\end{cousinweb}

\begin{grammarlens}[title={What you've built}]
One more word for saying what comes next.
\end{grammarlens}

\subsection*{Your turn}
\begin{itemize}
\raggedright
  \item \emph{Say it:} \emph{kaṭhpūtḷī}
  \item \emph{Say it:} \emph{kaṭhpūtḷī}, once more
  \item \emph{Say it:} \emph{kaṭhpūtḷī nāchai hai.} — The puppet dances.
  \item \emph{Recall:} say the Marwadi for a show, then the Marwadi for to wander about, then say \emph{kaṭhpūtḷī} again
\end{itemize}

\begin{tcolorbox}[breakable,colback=teal!4,colframe=teal!35!black,title={Before you move on}]
What does \mw{कठपूतळी} mean? (\textquotedblleft{}A puppet\textquotedblright{}.) Say it once more.
\end{tcolorbox}
\section[\texorpdfstring{\mw{लाडू} (lāḍū)}{लाडू (lāḍū)}]{\mw{लाडू} (lāḍū) — a laddu (a round sweet)}

\label{lesson:MW-C224-ladu}



\begin{quote}
Before the new one: say the Marwadi for to wander about, then the Marwadi for a puppet.
\end{quote}

\subsection*{You'll want to know: \mw{लाडू}}
\textbf{\mw{लाडू}} — \emph{lāḍū} — \textquotedblleft{}a laddu (a round sweet)\textquotedblright{}.

\begin{quote}
\mw{म्हैं} \mw{लाडू} \mw{खायो।}
\end{quote}

\begin{quote}
\emph{mhaĩ lāḍū khāyo.}
\end{quote}

\begin{quote}
I ate a laddu.
\end{quote}

A masculine noun.

\begin{cousinweb}
\textbf{\mw{लाडू}} (\emph{lāḍū}) is inherited from Sanskrit \emph{laḍḍu}, \textquotedblleft{}a kind of sweetmeat\textquotedblright{}; CDIAL lists Hindi and Gujarati \emph{lāḍū}, a ball of sweetened meal. CDIAL gives no Marwari form in this entry (CDIAL 10926).
\end{cousinweb}

\begin{grammarlens}[title={What you've built}]
One more word for saying what comes next.
\end{grammarlens}

\subsection*{Your turn}
\begin{itemize}
\raggedright
  \item \emph{Say it:} \emph{lāḍū}
  \item \emph{Say it:} \emph{lāḍū}, once more
  \item \emph{Say it:} \emph{mhaĩ lāḍū khāyo.} — I ate a laddu.
  \item \emph{Recall:} say the Marwadi for to wander about, then the Marwadi for a puppet, then say \emph{lāḍū} again
\end{itemize}

\begin{tcolorbox}[breakable,colback=teal!4,colframe=teal!35!black,title={Before you move on}]
What does \mw{लाडू} mean? (\textquotedblleft{}A laddu (a round sweet)\textquotedblright{}.) Say it once more.
\end{tcolorbox}
\section[\texorpdfstring{\mw{उछळणो} (uchhaḷṇo)}{उछळणो (uchhaḷṇo)}]{\mw{उछळणो} (uchhaḷṇo) — to leap up; to bounce}

\label{lesson:MW-C224-uchhalno}



\begin{quote}
Before the new one: say the Marwadi for a puppet, then the Marwadi for a laddu.
\end{quote}

\subsection*{You'll want to know: \mw{उछळणो}}
\textbf{\mw{उछळणो}} — \emph{uchhaḷṇo} — \textquotedblleft{}to leap up; to bounce\textquotedblright{}.

\begin{quote}
\mw{छोरो} \mw{ढोल} \mw{सुणै} \mw{है} \mw{अर} \mw{उछळै} \mw{है।}
\end{quote}

\begin{quote}
\emph{chhoro ḍhol suṇai hai ar uchhaḷai hai.}
\end{quote}

\begin{quote}
The boy hears the drum and leaps up.
\end{quote}

Past (he …): \mw{वो} \mw{छोरो} \mw{उछळ्यो}, \emph{vo chhoro uchhaḷyo}.

\begin{cousinweb}
\textbf{\mw{उछळणो}} (\emph{uchhaḷṇo}) is inherited from a reconstructed \emph{*ut-śalati}, \textquotedblleft{}springs up\textquotedblright{}; CDIAL lists the Old Marwari verb \emph{ūchaḷaï}, \textquotedblleft{}flies up, flutters\textquotedblright{} (CDIAL 1843).
\end{cousinweb}

\begin{grammarlens}[title={What you've built}]
One more word for saying what comes next.
\end{grammarlens}

\subsection*{Your turn}
\begin{itemize}
\raggedright
  \item \emph{Say it:} \emph{uchhaḷṇo}
  \item \emph{Say it:} \emph{uchhaḷṇo}, once more
  \item \emph{Say it:} \emph{chhoro ḍhol suṇai hai ar uchhaḷai hai.} — The boy hears the drum and leaps up.
  \item \emph{Recall:} say the Marwadi for a puppet, then the Marwadi for a laddu, then say \emph{uchhaḷṇo} again
\end{itemize}

\begin{tcolorbox}[breakable,colback=teal!4,colframe=teal!35!black,title={Before you move on}]
What does \mw{उछळणो} mean? (\textquotedblleft{}To leap up; to bounce\textquotedblright{}.) Say it once more.
\end{tcolorbox}
\section[\texorpdfstring{\mw{नगाड़ो} (nagāṛo)}{नगाड़ो (nagāṛo)}]{\mw{नगाड़ो} (nagāṛo) — a kettledrum}

\label{lesson:MW-C224-nagaro}



\begin{quote}
Before the new one: say the Marwadi for a laddu, then the Marwadi for to leap up.
\end{quote}

\subsection*{You'll want to know: \mw{नगाड़ो}}
\textbf{\mw{नगाड़ो}} — \emph{nagāṛo} — \textquotedblleft{}a kettledrum\textquotedblright{}.

\begin{quote}
\mw{म्हैं} \mw{नगाड़ो} \mw{सुण्यो।}
\end{quote}

\begin{quote}
\emph{mhaĩ nagāṛo suṇyo.}
\end{quote}

\begin{quote}
I heard the kettledrum.
\end{quote}

A masculine noun.

\begin{cousinweb}
\textbf{\mw{नगाड़ो}} (\emph{nagāṛo}) is usually traced to Persian \emph{naqqāra}, \textquotedblleft{}a kettledrum\textquotedblright{}, which Persian took from Arabic; Marwadi says \emph{g} for its \emph{q}.
\end{cousinweb}

\begin{grammarlens}[title={What you've built}]
One more word for saying what comes next.
\end{grammarlens}

\subsection*{Your turn}
\begin{itemize}
\raggedright
  \item \emph{Say it:} \emph{nagāṛo}
  \item \emph{Say it:} \emph{nagāṛo}, once more
  \item \emph{Say it:} \emph{mhaĩ nagāṛo suṇyo.} — I heard the kettledrum.
  \item \emph{Recall:} say the Marwadi for a laddu, then the Marwadi for to leap up, then say \emph{nagāṛo} again
\end{itemize}

\begin{tcolorbox}[breakable,colback=teal!4,colframe=teal!35!black,title={Before you move on}]
What does \mw{नगाड़ो} mean? (\textquotedblleft{}A kettledrum\textquotedblright{}.) Say it once more.
\end{tcolorbox}
